% GENERATED FILE. Edit canonical lessons, then run npm run generate:books.
% canonical-source-hash: fnv1a64:bde556f04fd7fdad
% canonical-lessons: PA-C334-unn, PA-C334-kattna-spin, PA-C334-bunna, PA-C334-julaha, PA-C334-rangna

\chapter{Wool and Cloth}
\label{ch:pa-a2-wool-and-cloth}
\hlchaptermodality{334}

\begin{chapteropening}
\textbf{By the end of this chapter:} \emph{I can say who spins, weaves and dyes.}

Complete the last lesson of chapter 334: I can say who spins, weaves and dyes.
\end{chapteropening}

\section[\texorpdfstring{\pa{ਉੱਨ} (unn)}{ਉੱਨ (unn)}]{\pa{ਉੱਨ} (unn) — wool}

\label{lesson:PA-C334-unn}



\begin{quote}
Before the new one: say the Punjabi for a cobbler, then the Punjabi for a gardener.
\end{quote}

\subsection*{You'll want to know: \pa{ਉੱਨ}}
\textbf{\pa{ਉੱਨ}} — \emph{unn} — \textquotedblleft{}wool\textquotedblright{}.

\begin{quote}
\pa{ਸਰਦੀ} \pa{ਲਈ} \pa{ਦਾਦੀ} \pa{ਉੱਨ} \pa{ਖ਼ਰੀਦਦੀ} \pa{ਹੈ।}
\end{quote}

\begin{quote}
\emph{sardī laī dādī unn kharīddī hai.}
\end{quote}

\begin{quote}
Grandmother buys wool for the winter.
\end{quote}

A feminine noun.

\begin{cousinweb}
\textbf{\pa{ਉੱਨ}} (\emph{unn}) is inherited from Sanskrit \emph{ūrṇā}, \textquotedblleft{}wool\textquotedblright{}; CDIAL lists Punjabi \emph{unn}, \textquotedblleft{}wool\textquotedblright{}, under it (CDIAL 2424).
\end{cousinweb}

\begin{grammarlens}[title={What you've built}]
One more word for this chapter.
\end{grammarlens}

\subsection*{Your turn}
\begin{itemize}
\raggedright
  \item \emph{Say it:} \emph{unn}
  \item \emph{Say it:} \emph{unn}, once more
  \item \emph{Say it:} \emph{sardī laī dādī unn kharīddī hai.} — Grandmother buys wool for the winter.
  \item \emph{Recall:} say the Punjabi for a cobbler, then the Punjabi for a gardener, then say \emph{unn} again
\end{itemize}

\begin{tcolorbox}[breakable,colback=teal!4,colframe=teal!35!black,title={Before you move on}]
What does \pa{ਉੱਨ} mean? (\textquotedblleft{}Wool\textquotedblright{}.) Say it once more.
\end{tcolorbox}
\section[\texorpdfstring{\pa{ਕੱਤਣਾ} (kattṇā)}{ਕੱਤਣਾ (kattṇā)}]{\pa{ਕੱਤਣਾ} (kattṇā) — to spin thread}

\label{lesson:PA-C334-kattna-spin}



\begin{quote}
Before the new one: say the Punjabi for a gardener, then the Punjabi for wool.
\end{quote}

\subsection*{You'll want to know: \pa{ਕੱਤਣਾ}}
\textbf{\pa{ਕੱਤਣਾ}} — \emph{kattṇā} — \textquotedblleft{}to spin thread\textquotedblright{}.

\begin{quote}
\pa{ਦਾਦੀ} \pa{ਸ਼ਾਮ} \pa{ਨੂੰ} \pa{ਉੱਨ} \pa{ਕੱਤਦੀ} \pa{ਹੈ।}
\end{quote}

\begin{quote}
\emph{dādī shām nū̃ unn kattdī hai.}
\end{quote}

\begin{quote}
Grandmother spins wool in the evening.
\end{quote}

Past (I …): \pa{ਮੈਂ} \pa{ਉੱਨ} \pa{ਕੱਤੀ}, \emph{maiṁ unn kattī}.

\begin{cousinweb}
\textbf{\pa{ਕੱਤਣਾ}} (\emph{kattṇā}) is inherited from a reconstructed \emph{*kartati}, \textquotedblleft{}spins\textquotedblright{}; CDIAL lists Punjabi \emph{kattṇā}, \textquotedblleft{}spins\textquotedblright{}, under it (CDIAL 2855).
\end{cousinweb}

\begin{grammarlens}[title={What you've built}]
One more word for this chapter.
\end{grammarlens}

\subsection*{Your turn}
\begin{itemize}
\raggedright
  \item \emph{Say it:} \emph{kattṇā}
  \item \emph{Say it:} \emph{kattṇā}, once more
  \item \emph{Say it:} \emph{dādī shām nū̃ unn kattdī hai.} — Grandmother spins wool in the evening.
  \item \emph{Recall:} say the Punjabi for a gardener, then the Punjabi for wool, then say \emph{kattṇā} again
  \item \emph{Write it:} copy \textbf{\pa{ਕੱਤਣਾ}} once, checking the whole word
\end{itemize}

\begin{tcolorbox}[breakable,colback=teal!4,colframe=teal!35!black,title={Before you move on}]
What does \pa{ਕੱਤਣਾ} mean? (\textquotedblleft{}To spin thread\textquotedblright{}.) Say it once more.
\end{tcolorbox}
\section[\texorpdfstring{\pa{ਬੁਣਨਾ} (buṇnā)}{ਬੁਣਨਾ (buṇnā)}]{\pa{ਬੁਣਨਾ} (buṇnā) — to knit}

\label{lesson:PA-C334-bunna}



\begin{quote}
Before the new one: say the Punjabi for wool, then the Punjabi for to spin thread.
\end{quote}

\subsection*{You'll want to know: \pa{ਬੁਣਨਾ}}
\textbf{\pa{ਬੁਣਨਾ}} — \emph{buṇnā} — \textquotedblleft{}to knit\textquotedblright{}.

\begin{quote}
\pa{ਮਾਤਾ} \pa{ਜੀ} \pa{ਮੇਰੇ} \pa{ਲਈ} \pa{ਸਵੈਟਰ} \pa{ਬੁਣਦੇ} \pa{ਹਨ।}
\end{quote}

\begin{quote}
\emph{mātā jī mere laī savaiṭar buṇde han.}
\end{quote}

\begin{quote}
Mother knits a sweater for me.
\end{quote}

Past (I …): \pa{ਮੈਂ} \pa{ਇੱਕ} \pa{ਸਵੈਟਰ} \pa{ਬੁਣਿਆ}, \emph{maiṁ ikk savaiṭar buṇiā}.

\begin{cousinweb}
\textbf{\pa{ਬੁਣਨਾ}} (\emph{buṇnā}) is inherited from a reconstructed \emph{*vināti}, \textquotedblleft{}weaves\textquotedblright{}; CDIAL lists Punjabi \emph{bunṇā}, \textquotedblleft{}to weave\textquotedblright{}, under it (CDIAL 11773).
\end{cousinweb}

\begin{grammarlens}[title={What you've built}]
One more word for this chapter.
\end{grammarlens}

\subsection*{Your turn}
\begin{itemize}
\raggedright
  \item \emph{Say it:} \emph{buṇnā}
  \item \emph{Say it:} \emph{buṇnā}, once more
  \item \emph{Say it:} \emph{mātā jī mere laī savaiṭar buṇde han.} — Mother knits a sweater for me.
  \item \emph{Recall:} say the Punjabi for wool, then the Punjabi for to spin thread, then say \emph{buṇnā} again
  \item \emph{Write it:} copy \textbf{\pa{ਬੁਣਨਾ}} once, checking the whole word
\end{itemize}

\begin{tcolorbox}[breakable,colback=teal!4,colframe=teal!35!black,title={Before you move on}]
What does \pa{ਬੁਣਨਾ} mean? (\textquotedblleft{}To knit\textquotedblright{}.) Say it once more.
\end{tcolorbox}
\section[\texorpdfstring{\pa{ਜੁਲਾਹਾ} (julāhā)}{ਜੁਲਾਹਾ (julāhā)}]{\pa{ਜੁਲਾਹਾ} (julāhā) — a weaver}

\label{lesson:PA-C334-julaha}



\begin{quote}
Before the new one: say the Punjabi for to spin thread, then the Punjabi for to knit.
\end{quote}

\subsection*{You'll want to know: \pa{ਜੁਲਾਹਾ}}
\textbf{\pa{ਜੁਲਾਹਾ}} — \emph{julāhā} — \textquotedblleft{}a weaver\textquotedblright{}.

\begin{quote}
\pa{ਜੁਲਾਹਾ} \pa{ਸਾਰਾ} \pa{ਦਿਨ} \pa{ਖੇਸ} \pa{ਬੁਣਦਾ} \pa{ਹੈ।}
\end{quote}

\begin{quote}
\emph{julāhā sārā din khes buṇdā hai.}
\end{quote}

\begin{quote}
The weaver weaves thick cotton sheets all day.
\end{quote}

A masculine noun.

\begin{cousinweb}
\textbf{\pa{ਜੁਲਾਹਾ}} (\emph{julāhā}) is a loan, usually traced to Persian \emph{jūlāh}, \textquotedblleft{}a weaver\textquotedblright{}.
\end{cousinweb}

\begin{grammarlens}[title={What you've built}]
One more word for this chapter.
\end{grammarlens}

\subsection*{Your turn}
\begin{itemize}
\raggedright
  \item \emph{Say it:} \emph{julāhā}
  \item \emph{Say it:} \emph{julāhā}, once more
  \item \emph{Say it:} \emph{julāhā sārā din khes buṇdā hai.} — The weaver weaves thick cotton sheets all day.
  \item \emph{Recall:} say the Punjabi for to spin thread, then the Punjabi for to knit, then say \emph{julāhā} again
  \item \emph{Write it:} copy \textbf{\pa{ਜੁਲਾਹਾ}} once, checking the whole word
\end{itemize}

\begin{tcolorbox}[breakable,colback=teal!4,colframe=teal!35!black,title={Before you move on}]
What does \pa{ਜੁਲਾਹਾ} mean? (\textquotedblleft{}A weaver\textquotedblright{}.) Say it once more.
\end{tcolorbox}
\section[\texorpdfstring{\pa{ਰੰਗਣਾ} (raṅgṇā)}{ਰੰਗਣਾ (raṅgṇā)}]{\pa{ਰੰਗਣਾ} (raṅgṇā) — to dye}

\label{lesson:PA-C334-rangna}



\begin{quote}
Before the new one: say the Punjabi for to knit, then the Punjabi for a weaver.
\end{quote}

\subsection*{You'll want to know: \pa{ਰੰਗਣਾ}}
\textbf{\pa{ਰੰਗਣਾ}} — \emph{raṅgṇā} — \textquotedblleft{}to dye\textquotedblright{}.

\begin{quote}
\pa{ਦਾਦੀ} \pa{ਚਿੱਟੀ} \pa{ਉੱਨ} \pa{ਨੂੰ} \pa{ਲਾਲ} \pa{ਰੰਗਦੀ} \pa{ਹੈ।}
\end{quote}

\begin{quote}
\emph{dādī chiṭṭī unn nū̃ lāl raṅgdī hai.}
\end{quote}

\begin{quote}
Grandmother dyes the white wool red.
\end{quote}

Past (I …): \pa{ਮੈਂ} \pa{ਆਪਣੀ} \pa{ਚੁੰਨੀ} \pa{ਪੀਲੀ} \pa{ਰੰਗੀ}, \emph{maiṁ āpṇī chunnī pīlī raṅgī}.

\begin{cousinweb}
\textbf{\pa{ਰੰਗਣਾ}} (\emph{raṅgṇā}) is inherited from a reconstructed \emph{*raṅgayati}, \textquotedblleft{}dyes\textquotedblright{}; CDIAL lists Punjabi \emph{raṅgṇā}, \textquotedblleft{}to dye, paint\textquotedblright{}, under it (CDIAL 10570).
\end{cousinweb}

\begin{grammarlens}[title={What you've built}]
One more word for this chapter.
\end{grammarlens}

\subsection*{Your turn}
\begin{itemize}
\raggedright
  \item \emph{Say it:} \emph{raṅgṇā}
  \item \emph{Say it:} \emph{raṅgṇā}, once more
  \item \emph{Say it:} \emph{dādī chiṭṭī unn nū̃ lāl raṅgdī hai.} — Grandmother dyes the white wool red.
  \item \emph{Recall:} say the Punjabi for to knit, then the Punjabi for a weaver, then say \emph{raṅgṇā} again
  \item \emph{Write it:} copy \textbf{\pa{ਰੰਗਣਾ}} once, checking the whole word
\end{itemize}

\begin{tcolorbox}[breakable,colback=teal!4,colframe=teal!35!black,title={Before you move on}]
What does \pa{ਰੰਗਣਾ} mean? (\textquotedblleft{}To dye\textquotedblright{}.) Say it once more.
\end{tcolorbox}
